\section{模型的稳健性与灵敏度分析}

\subsection{实验设计原则}

四个问题的稳健性、灵敏度与消融分析遵循三条统一原则。第一，\textbf{检验标准事先设定}：核心结论、检验标准、阈值、扰动安排与样本量在计算前固定，事后不得修改；凡未能通过或被证伪的项一律如实列出，不通过调参、放宽阈值或挑选指标使其“通过”。第二，\textbf{不污染主方案}：随机规划、鲁棒优化、预测器替换、决策加密、前瞻深度等对照全部只进消融，主结果始终出自“确定性线性规划 + 题面规定的信息结构”。第三，\textbf{基线门槛}：每批扰动实验必须首先逐位复现主结果，否则整批作废，以保证对照之间的可比性——四问的基线复现相对差全部为 0（或不超过 $10^{-14}$）。

\subsection{稳定性分级汇总}

四问的扰动实验规模与结论汇总于表~\ref{tab:rob}。需要强调，分级所用的检验标准集合在计算前即已固定，因此“未通过”是结论而不是失败：它给出的是模型结论的适用边界。

\begin{table}[htbp]
\centering
\footnotesize
\caption{四个问题的扰动实验规模、分级与主要敏感源}
\label{tab:rob}
\begin{tabular}{p{1.5cm}rp{3.0cm}p{2.2cm}p{4.9cm}}
\toprule
问题 & 情景数 & 未通过检验标准 & 分级 & 主要敏感源与关键量 \\
\midrule
问题一 & 948（947 可行） & 无 & 稳定 & 往返效率 $\eta$（$0.9\to0.8$ 时 $+8.885\%$）；初值影响 $\le0.07\%$ \\
问题二 & 149 次全年 LP & 2 条（已归因） & 脆弱（需缩小适用范围） & $\eta$（$\pm10\%$ 使年度总量约 $\pm7\%$、单日量约 $\pm20\%$） \\
问题三 & 156 情景 & 1 条（噪声跨档均值） & 条件稳定（需给出适用边界） & $\eta$（单侧最大 10.9715\%）；初值对交付期无影响 \\
链 \emph{4-2} & 162 & 无 & 稳定 & 往返效率（$-20\%$ 档 $+13.345\%$）与储电量上限（$-20\%$ 档 $+5.141\%$） \\
链 \emph{4-3} & 198（195 成功） & 1 条（$\sigma=10\%$ 噪声族） & 条件稳定（需给出适用边界） & 光伏预报偏差（$+10\%$ 档 $+23.25\%$）与光伏噪声 \\
\bottomrule
\end{tabular}
\end{table}

\begin{figure}[htbp]
\centering
\includegraphics[width=0.490\textwidth]{fig19_robust_tornado.png}
\caption{参数敏感性的龙卷风图（问题三；其余问题为同类图）}
\label{fig:rob1}
\end{figure}

\subsection{敏感性主因的识别}

四个问题在敏感源上给出一致而清晰的图景，其中三点具有一般性。

第一，\textbf{往返效率是唯一一阶参数}。在四个问题上效率的影响都远大于其他参数：问题一 $\eta$ 由 0.9 降至 0.8 时费用上升 8.885\%、升至 1.0 时下降 7.828\%；问题二 $\eta$ 取 0.72 时费用上升 14.078\%，弹性约 $-0.66$ 至 $-0.70$；问题三单侧最大影响 10.9715\%；链 \emph{4-2} 为 13.345\%。相反，初始储电量的影响在所有问题上都可忽略（问题一不超过 0.07\%，问题二不超过 0.0042\%，问题三对交付期费用完全没有影响，链 \emph{4-2} 无响应）。这一对比有明确的物理解释：效率直接决定储能搬运电能的成本，而初值只影响开局的少量调节裕度；文献亦指出实际效率随荷电状态与功率变化，恒效率属工程简化，因此该参数既是最敏感项、也是最需要实测标定的一项\cite{ref-vykhodtsev-2022-bess-review,ref-huang-2022-dynamic-efficiency}。

第二，\textbf{敏感性的统计方式依赖必须随数字一起说明}。问题二的交付期费用稳定性检验之所以未通过，正是因为往返效率在日尺度量与年度总量上的敏感度相差约三倍（数值见问题二一节），因此本文在引用表~\ref{tab:p2a}、表~\ref{tab:p2b} 的单日数值时一律附带该声明。

第三，\textbf{“价格波动是第一不确定性源”按最大值统计时不成立}。这是本文如实列出的一个反例：按平均值统计时价格类扰动的影响确实居前，但按最大绝对值统计时价格类扰动的最大影响为 10.0\%，低于效率扰动的 13.345\%（链 \emph{4-2}）与 11.700\%（链 \emph{4-3}）。真正的一阶通道在链 \emph{4-2} 上是题面硬参数往返效率，在链 \emph{4-3} 上是光伏预报误差。本文不使用“价格是第一不确定性源”这一未经审慎表述的结论。

此外，链 \emph{4-3} 的未通过检验标准给出了本问题最有价值的适用边界：光伏与负载的白噪声在 $\sigma=5\%$ 下使费用上升约 3.9\%、在 $\sigma=10\%$ 下上升 11.2\%，而价格侧同类扰动影响明显更小。这说明在“预报驱动的三层决策 + 单侧紧急购电 + 非对称双向结算”这一结构下，\textbf{模型对光伏侧误差的暴露远大于对价格侧误差的暴露}，机制是紧急购电项的凸性把误差系统性地转化为费用。对工程实践的直接含义是：若要降低总费用，优先级应是提高光伏预报精度，而不是提高价格预报精度。

\subsection{模型族对比的整体结论}

跨四个问题的消融给出四条一致结论。

第一，\textbf{线性松弛无损，无需二元变量}。问题一中加入 144 个互补二元变量的整数模型与线性模型目标值完全相同；问题二中加入 52\,560 个二元变量的全规模整数模型只比线性模型高 4.355178 元（相对差 $3.2\times10^{-7}$），代价是求解时间从 5.8 s 增到 301.6 s 且未在时限内证明最优。引理 1 给出了这一现象的解析解释，也说明在规模较大时不必按混合整数规划建模。

第二，\textbf{独立第二实现与精确分解都逐位复现主结果}。问题二与问题三的独立第二实现以及问题二的精确逐日分解，都给出与主模型逐位相同的费用（做法与数值见问题二、问题三的模型族对照），这排除了主结果为实现假象的可能。

第三，\textbf{建模选择的影响可以大于机制价值}。问题三中把光伏预报的降尺度从整点锚定线性插值改为分段常数前向保持，会使费用上升 15.83\%（2\,561\,954.06 元），远大于调整机制本身 2.02\% 的价值；问题四中价格预测器的选择对费用影响不超过 0.28\%。这提示\textbf{数据处理方式的选择往往比模型结构的精细化更值得审慎对待}。

第四，\textbf{时域安排具有可量化的经济含义}。跨日状态连续、终端条件与日周期约束三项分别对应 26\,894.74 元、2217.08 元与 17\,135.02 元的费用差异；而“日边界不携带电量”这一在四条链上都出现的现象，被滚动时域对照证明只是“逐日决策加终端自由”结构的产物——把前瞻深度提高到 3 天以上，日末储电量即整体抬离下限（图~\ref{fig:rob2}，逐项数值见问题四的消融对照）。

\begin{figure}[htbp]
\centering
\includegraphics[width=0.535\textwidth]{fig20_rh_depth.png}
\caption{前瞻深度对费用的影响与日边界储电量行为：日边界不携带电量只在特定结构下成立}
\label{fig:rob2}
\end{figure}

\subsection{如实列出的反例与未通过项}

按“不得强凑单调性”的原则，本文记录以下六项与预期不符或未通过的结果：

\begin{enumerate}[leftmargin=2.2em,itemsep=1pt]
  \item 问题二的可行率检验未通过：唯一未通过的样本的目标值与基线逐位相同、解可行，未通过项是“返回解是否满足充放电互补性”，属退化最优面上的返回解选择问题。
  \item 问题二的交付表稳定检验未通过：$\eta=0.99$ 使 2025 年 6 月 21 日全天购电量下降 19.10\%，属有利方向的变化，但因检验标准取绝对值而被判失败；反向（$\eta=0.81$）最多使单日购电费上升 12.97\%。
  \item 问题三的噪声检验未通过：跨噪声幅度档的族间均值差 7.29\% 超过 3\% 阈值，而同幅度档内仅 0.06\%，属把两类统计量绑在同一检验标准里的设计问题。
  \item 链 \emph{4-3} 的噪声检验未通过：$\sigma=10\%$ 族均值上升 10.5919\% 超过 5\% 阈值，机制为紧急购电项的凸性，属结构性而非实现缺陷。
  \item 决策加密无收益：“每 3 小时”与“每 1 小时”的决策密度使费用分别上升 398.90 元与 590.07 元。本文不将其解读为“越密越差”，只记录“在题面给定的四个发布时刻下无收益”。
  \item 组合对照不叠加收益：两阶段随机与前瞻深度的组合方案在同一窗口内的期望费用比点预测方案高 2.6959\%（算例设定见问题四一节）。该结论只对 12 个代表日的窗口成立，本文不作外推。
\end{enumerate}

\subsection{对“是否需要更多预报、是否需要随机优化”的回答}

题面在问题三中提出了“是否需要引入其他时刻的预报制定调整购电策略”这一问题，本文把它拆成三个递进层面回答，并以计算成本作为约束；各层面的完整证据分别见问题三与问题四，此处给出汇总与取舍判断。

\textbf{层面一：要不要有调整？} 要。加入 6:00、12:00、18:00 三个调整时刻后的净收益为正（2.02\%），机制是减少 5 倍价紧急购电、代价是新增偏差违约金，两项金额已在问题三的费用分解中给出。

\textbf{层面二：在四个时刻之上是否还要加密？} 不需要。加密决策时刻并不增加信息量，只是把同一套预报用得更频繁，费用不降反升（金额见上一小节反例第 5 条）。

\textbf{层面三：要不要为了“看更远”而引入滚动优化？} 有收益但边际递减：两链交付期费用在 3 天前瞻上分别再降约 0.18\% 与 0.15\%，提高到 7 天与 14 天几乎不再改善。因此真正的价值来自“至少看 3 天”，而不是“看得越远越好”；同时该改动会使日末储电量脱离下限，从根本上改变模型的跨日行为。

\textbf{关于随机优化}：本文构造的两阶段随机规划对照在链 \emph{4-2} 上使同窗期望费用降低 3.2860\%，说明价格不确定性确实存在可对冲的部分；但链 \emph{4-3} 的随机对照因事先设定的自由追索而退化，其 19.02\% 的“降幅”不可作为对冲收益。综合收益幅度（不超过 3.3\%）、场景建模所需的历史样本量与额外计算预算，本文维持“价格点预测 + 确定性线性规划”为主方案，并明确指出：若要改换主方案，必须先补足价格与光伏误差的实测分布证据\cite{ref-akbari-2019-stochastic-robust-arbitrage,ref-oconnor-2025-conformal-price-forecast}。
